\chapter{Computational Framework for the Average Metrics}
\label{app:computational_framework}

This appendix sets out the computational framework behind the average metrics, which extends the \texttt{RBesT} package \citep{weber_applying_2021}. Decision rules are built with \texttt{decision2S}, the decision boundary of Equation~\eqref{eq:app_boundary} with \texttt{decision2S\_boundary}, the conditional operating characteristics of Equation~\eqref{eq:oc2S_normal} with \texttt{oc2S}, and the averaging integral of Equation~\eqref{eq:avg_integral} with \texttt{integrate\_density\_log}.

Let $\bar{y}_T$ and $\bar{y}_C$ denote the sample means for the treatment and control arms, respectively.

Success is determined by a decision function $D(\bar{y}_T, \bar{y}_C) \in \{0, 1\}$, which evaluates whether the posterior distribution of the treatment contrast, 
  $\delta = \theta_C - \theta_T$, satisfies a set of predefined conditions. For instance, a standard design requires the posterior probability to 
  exceed a specific threshold: $\Pr(\delta > q \mid \bar{y}_T, \bar{y}_C) > p$ for threshold $q$ and probability $p$.

When evaluating multiple conditions simultaneously, such as in \textit{dual criterion} designs requiring both statistical significance and clinical relevance 
\citep{roychoudhury_beyond_2018}, the decision is formulated as the product of Heaviside step functions for each condition\footnote{Note that this product 
formulation naturally generalizes to an arbitrary number of conditions} $H(x)$, where $H(x) = 1$ for $x > 0$ and $0$ for $x \le 0$, and where $q_{c,i}$ and $p_{c,i}$ are the reference value and the posterior probability threshold of the $i$th criterion:

$$D(\bar{y}_T, \bar{y}_C) = \prod_{i} H\Big(\Pr(\delta \ge q_{c,i} \mid \bar{y}_T, \bar{y}_C) - p_{c,i}\Big)$$



The specification of priors, the sample sizes ($n_T$, $n_C$), and the decision function $D(\bar{y}_T, \bar{y}_C)$ uniquely define a conditional decision boundary 
$D_T(\bar{y}_C)$ for each observed control outcome $\bar{y}_C$,
\begin{equation} \label{eq:app_boundary}
D_T(\bar{y}_C) = \sup\{\bar{y}_T : D(\bar{y}_T, \bar{y}_C) = 1\}
\end{equation}
This is the largest value of $\bar{y}_T$ that still yields success, so success occurs when $\bar{y}_T$ falls below it.

Given this boundary, we can compute the conditional probability of success $\Pr(\text{Success} \mid \theta_C, \theta_T)$ required in Equation~\eqref{eq:m_joint}. Assuming
 true parameters $\theta_C$ and $\theta_T$, the overall probability of a success is the expectation of crossing the boundary $D_T(\bar{y}_C)$ over the sampling
 distribution of the control data:
\begin{equation}
  \label{eq:oc2S_general}
  \Pr(\text{Success} \mid \theta_C, \theta_T) = \int_{-\infty}^{\infty} F_T\left(D_T(\bar{y}_C) \mid \theta_T\right) f_C\left(\bar{y}_C \mid \theta_C\right) \, d\bar{y}_C,
\end{equation}
where $F_T\left(\cdot \mid \theta_T\right)$ is the cumulative distribution function of the treatment statistic evaluated at the decision boundary, and
 $f_C\left(\cdot \mid \theta_C\right)$ is the probability density function of the control statistic.

Because we are working with normally distributed endpoints, the sampling distributions for the sample means are Normal: 
$\bar{y}_m \sim \mathcal{N}(\theta_m, v_m^2)$ for $m \in \{T, C\}$, where $v_m = \sigma_m / \sqrt{n_m}$. Consequently, the general integral in 
Equation~\eqref{eq:oc2S_general} can be evaluated directly using the standard Normal CDF ($\Phi$) and PDF ($\phi$):
\begin{equation}
  \label{eq:oc2S_normal}
  \Pr(\text{Success} \mid \theta_C, \theta_T) = \int_{-\infty}^{\infty} \Phi\left(\frac{D_T(\bar{y}_C) - \theta_T}{v_T}\right) \frac{1}{v_C} \phi\left(\frac{\bar{y}_C - \theta_C}{v_C}\right) \, d\bar{y}_C.
\end{equation}
 

\section{The Decision Boundary under a Mixture Analysis Prior}

Equation~\eqref{eq:oc2S_normal} presumes the boundary $D_T(\bar{y}_C)$ of
Equation~\eqref{eq:app_boundary} is available. For a single normal analysis prior it is available in closed form: the
posterior mean of the contrast is affine in $\bar{y}_C$, the posterior variance does not depend on
the data, and the success condition reduces to the fixed threshold of
Equation~\eqref{eq:threshold_k}, so $D_T$ is an affine function of $\bar{y}_C$ with slope $W$.

A mixture analysis prior breaks the last two of these. The posterior is again a normal mixture, but its
component weights are updated by $\bar{y}_C$, so the posterior variance depends on the data and the
success condition can no longer be inverted analytically. The boundary is instead located
numerically, by solving $\Pr(\delta > q \mid \bar{y}_T, \bar{y}_C) = p$ for $\bar{y}_T$ at each
$\bar{y}_C$, and it is no longer affine: for the three component MAP prior of
Section~\ref{subsec:formulating_priors} it is visibly curved over the plausible range of
$\bar{y}_C$, whereas for the single normal its curvature is numerically zero.

This is the computational reason the mixture designs of Chapter~\ref{sec:Applications} have no
closed form, and it is a property of the analysis prior alone. The design prior enters only as the
integrating measure in Equation~\eqref{eq:avg_integral}, so a mixture design prior costs nothing
beyond the numerical integration that is required in any case.

\section{Averaging over the Design Prior}

Equation~\eqref{eq:oc2S_normal} gives the probability of success at a fixed pair $(\theta_C, \theta_T)$.
The average metrics of Section~\ref{subsec:metrics_bayesian} replace that fixed pair by a draw from the
design prior. Writing $\theta_T = \theta_C - \delta$, so that the treatment effect is held at the value of
interest while the control parameter varies, the average probability of success is
\begin{equation} \label{eq:avg_integral}
  \Pr(\text{Success} \mid \delta, \pi_D)
  = \int_{-\infty}^{\infty} \Pr\!\left(\text{Success} \given \theta_C,\, \theta_C - \delta\right)
    \pi_D(\theta_C)\, d\theta_C .
\end{equation}
Setting $\delta = 0$ returns the avgT1E and any other $\delta$ returns the avgPower at that effect, so a
single integral produces both metrics. The integrand is the conditional curve of
Equation~\eqref{eq:oc2S_normal}, which is bounded in $[0,1]$ and smooth in $\theta_C$, and the integral is
evaluated numerically on the logit probability scale over the central mass of $\pi_D$, discarding a mass
$\varepsilon/2$ in each tail, with $\varepsilon = 10^{-6}$ throughout. For a Normal mixture design prior no
simulation is required at any point. The routine used, \texttt{integrate\_density\_log}, is internal to
\texttt{RBesT} rather than part of its exported interface, so this step depends on an implementation detail
of that package.

With avgT1E and avgPower available from Equation~\eqref{eq:avg_integral} at $\delta = 0$ and at the effect of
interest respectively, the fixed design EUII of Equation~\eqref{eq:euii_avg} requires no further integration:
the total sample size $n_{\text{trial}}$ is deterministic, so obtaining the EUII is a single division on the
log scale rather than an average over $N$.

\section{Monte Carlo Evaluation of Sequential Designs}

A group sequential design has no closed form counterpart to Equation~\eqref{eq:avg_integral}, because success
at look $k$ depends on the whole history of the trial. The operating characteristics are therefore obtained by
simulation, with the three features below keeping the cost manageable.

\paragraph{Boundaries are precomputed once.} The decision boundary at each look is a property of the analysis,
not of the truth: it is determined by the priors, the sample sizes and the decision rule, and depends neither on
$\theta_C$ nor on $\delta$. Each look therefore contributes a function $D_{T,k}(\bar{y}_{C,k})$, built once by the averaging routine
before the simulation starts and reused for every replicate and every value of $\delta$.

\paragraph{Trials are simulated jointly across replicates and retired as they stop.} For $n_{\text{sim}}$ replicates the
control parameter is drawn from the design prior, $\theta_C^{(i)} \sim \pi_D$, and the treatment parameter is
set to $\theta_T^{(i)} = \theta_C^{(i)} - \delta$. At each look only the increment of new patients is
generated, and the running arm means are updated cumulatively. A replicate that crosses the efficacy boundary is
recorded as an efficacy stop at that look, one that crosses the futility boundary as a futility stop, and either
way it is removed from the active set, so later looks are simulated only for the trials still running. The
decision rule at each look may test several conditions at once, exactly as in the dual criterion case of
Equation~\eqref{eq:app_boundary}: \texttt{decision2S} accepts a vector of thresholds for this purpose, and the
boundary crossing check above applies unchanged, so no separate code path is needed for the DC.

\paragraph{The conditional case is the same computation.} Drawing every $\theta_C^{(i)}$ from $\pi_D$ gives
the average metrics, while fixing $\theta_C^{(i)} = \theta_C^*$ for all $i$ gives the conditional ones. The
conditional operating characteristics are thus the special case of a degenerate design prior, and no separate
code path is needed.

The stopping probabilities of Section~\ref{subsec:sequential} are the proportions of replicates retired at each
look, and summing the efficacy proportions over looks gives avgPower, or avgT1E when $\delta = 0$. A
replicate still running at the final look is recorded as a failure, since the final analysis admits only
success or failure (Section~\ref{sec:trial_setup}).

The same per look proportions supply the sequential EUII of Equation~\eqref{eq:seq_euii}. Running the simulation
once at $\delta = 0$ and once at the effect of interest gives, for each run, the per look efficacy proportion
divided by that run's own avgT1E or avgPower, which is exactly $\Pr(\text{stop at } k \mid \text{success}, H_i)$
of Equation~\eqref{eq:seq_inv_enp} for $H_i = H_0$ or $H_1$ respectively. The analogous ratio to
$1 - \text{avgT1E}$ or $1 - \text{avgPower}$ gives $\Pr(\text{stop at } k \mid \text{no success}, H_i)$ of
Equation~\eqref{eq:seq_inv_enm}. These two simulations therefore supply every term needed for
$\Erw[1/N_i^+]$ and $\Erw[1/N_i^-]$. Combining them into $\Erw[1/N^+]$ and $\Erw[1/N^-]$ through the
posterior probabilities of $H_0$ and $H_1$, and then into the EUII itself, is the closed form Bayes' rule
step of Section~\ref{subsubsec:seq_euii} and needs no further simulation.

\paragraph{Precision.} avgT1E and avgPower are each an average of Bernoulli stopping indicators over
$n_{\text{sim}}$ replicates, so each carries a Monte Carlo standard error of order
$\sqrt{\hat{p}(1-\hat{p})/n_{\text{sim}}}$. The EUII is a function of these two proportions and of the per
look proportions behind $\Erw[1/N^+]$ and $\Erw[1/N^-]$, and inherits their uncertainty through a nonlinear
transform, so it carries a larger error. At the replicate counts used here both are small: at
$n_{\text{sim}} = 2 \times 10^6$ an avgT1E near $0.05$ has a standard error below $2 \times 10^{-4}$, well
inside the three decimals reported, and the EUII is resolved to the leading digits shown. Every calibrated
design is nonetheless reevaluated once more, at a fresh seed and at twice the number of replicates used during
the calibration search, before it is reported: $n_{\text{sim}} = 10^6$ for the search and
$n_{\text{sim}} = 2 \times 10^6$ for this verification, in the calibrated comparisons of
Section~\ref{sec:calibrated_comparison}.
